% Generated draft: requires mathematical and source review.
\documentclass[11pt]{article}
\usepackage[margin=25mm]{geometry}
\usepackage{fontspec}
\setmainfont{Latin Modern Roman}
\usepackage{amsmath,amssymb,mathtools,unicode-math}
\setmathfont{Latin Modern Math}
\usepackage{enumitem,needspace,xurl,hyperref}
\setlength{\parindent}{0pt}
\setlength{\parskip}{5pt}
\setlength{\emergencystretch}{3em}
\newcommand{\N}{\mathbb{N}}
\newcommand{\Z}{\mathbb{Z}}
\newcommand{\Q}{\mathbb{Q}}
\newcommand{\R}{\mathbb{R}}
\newcommand{\C}{\mathbb{C}}
\newcommand{\F}{\mathbb{F}}
\newcommand{\A}{\mathbb{A}}
\newcommand{\PP}{\mathbb{P}}
\newcommand{\E}{\mathbb{E}}
\newcommand{\eps}{\varepsilon}
\newcommand{\dd}{\,\mathrm{d}}
\begin{document}
\section*{UnsolvedMath ID: 30005971}
\textbf{Source problem code:} catalogue entry
\subsection{Definitions and mathematical statement}
Let \(p\) denote a prime. For a nonzero rational number \(x\), write \(x=p^k a/b\), where \(k\in\Z\) and \(a,b\in\Z\) are not divisible by \(p\), and define \(v_p(x)=k\).

The supplied record poses three related questions. It does not specify the automorphic setting, the spaces \(M(\mathcal O_j)\), or the regulator conventions, so the following is a statement of the requested constructions and predictions with these qualifications left explicit.

\begin{enumerate}
\item Can one construct and characterize mod-\(p\) automorphic induction without passing through mod-\(p\) Galois representations? The fields, field extension, reductive groups, coefficient field, classes of automorphic objects, and properties intended to characterize induction are not specified in the supplied record.
\item For the spaces denoted by \(M(\mathcal O_1)\) and \(M(\mathcal O_2)\), can one predict the \(p\)-adic valuations of the relevant regulator ratios and, with the source qualification ``when it is nonzero,'' of
\[
\frac{\#H^i(M(\mathcal O_1),\Z)}{\#H^i(M(\mathcal O_2),\Z)}?
\]
Here \(H^i(-,\Z)\) denotes integral cohomology and \(\#\) denotes cardinality. To apply the rational valuation defined above to this displayed quotient, both cohomology groups must be finite. That finiteness is not supplied as a hypothesis. Replacing these groups by their torsion subgroups would change the statement and is not done here. If the groups are finite, their orders are positive, so the meaning of the source's additional nonzero qualification remains unclear.

The original statement renders the regulator expression as
\[
\frac{\operatorname{Reg}_i(M(\mathcal O_1))^2}{\operatorname{Reg}_i(M(\mathcal O_1))^2}.
\]
Taken literally, this quotient is \(1\) whenever defined. The cleaned statement instead refers to unspecified ``relevant ratios of regulators.'' The record does not establish whether the repeated \(\mathcal O_1\) is intentional, whether a denominator involving \(\mathcal O_2\) was intended, or whether the squares are essential. Consequently no corrected regulator quotient is asserted here.
\item Can one define \(p\)-adic regulators whose ratios equal the corresponding rational ratios of real regulators? More explicitly, for each intended pair of real regulators with nonzero denominator and rational quotient \(r\in\Q\), the requested compatibility is equality of the corresponding \(p\)-adic regulator quotient with the image of \(r\) under the canonical embedding \(\Q\hookrightarrow\Q_p\), provided the proposed regulators take values in \(\Q_p\). The actual target of a \(p\)-adic regulator, any necessary comparison map, the real regulator normalization, the relevant pairs, and whether regulators or their squares are compared are not specified.
\end{enumerate}

The permissible primes \(p\), degrees \(i\), and pairs \((\mathcal O_1,\mathcal O_2)\), together with their geometric and arithmetic hypotheses, are absent. No universal quantification over arbitrary orders or spaces is imposed. These omissions prevent a fully specified mathematical formulation from being recovered from this record alone.
\subsection{Short English statement}
The source asks for a construction and characterization of mod-\(p\) automorphic induction that avoids mod-\(p\) Galois representations, predictions for valuations of regulator and integral cohomology order ratios, and compatible \(p\)-adic regulators. The setting, finiteness assumptions, and regulator conventions are missing, and the original regulator formula repeats the same space in numerator and denominator. This statement-only draft preserves those unresolved issues and requires human mathematical review.
\subsection{Sources}
\begin{itemize}
\item Local registry: \texttt{data/unsolvedmath/problems.json}, supplied record with permanent UnsolvedMath ID \(30005971\), title ``Mod-p Automorphic Induction and p-Adic Regulators,'' source problem code \texttt{catalogue entry}. The supplied record gives proposed year \(2024\) and identifies no individual proposer.
\item Supplied source citation: \emph{Explicit Methods in Number Theory} (2025), Oberwolfach Reports, DOI \texttt{10.4171/owr/2024/40}, pp.\ 2309--2368. \url{https://doi.org/10.4171/owr/2024/40}. The background describes Oberwolfach Reports as documenting mathematical workshops and their open-problem discussions. The citation year \(2025\) is retained alongside the record's \(2024\) proposed year and year tag; their relationship has not been independently checked.
\item Supplied literature source: ``Bartel--Page, Vigneras orbifolds.'' \url{https://arxiv.org/abs/2407.07240}. This attribution is retained as supplied; the record provides no theorem statement or precise result reference establishing the intended formulation.
\item Literature assessment date: \texttt{2026-08-22}. This is supplied dated curation, not an independently verified assessment. The supplied assessment labels the questions open and reports that existing work gives exclusions and conditional mod-\(p\) results, without a general solution. The background expressly requests independent verification of time-sensitive status and cited claims. No such verification is asserted here.
\item The supplied registry labels the statement \texttt{corrected\_verified} and says that only formula delimiters, notation, grammar, and typography were repaired, with mathematical content and scope preserved. This is source metadata, not a mathematical verification of this draft; the regulator formula discrepancy remains unresolved.
\end{itemize}
\end{document}
